
\subsection{算子 $\mathfrak p$ 的代数表述 II：
$\mathfrak p$ 诱导 $L_\infty$ 模同态}
\label{subsec:7-4-11}

设 $L$ 是辛流形 $M$ 的相对 spin 拉格朗日子流形。
如第~\ref{subsec:7-4-3} 小节一样，
把 $H(M;\Lambda_{0,\mathrm{nov}})[1]$
视为平凡 $L_\infty$ 代数。
再考虑由次数分别为 $1$、$0$ 的两个元素
$e_1,e_0$ 生成的自由 $\Lambda_{0,\mathrm{nov}}$-模
\[
\Lambda_{0,\mathrm{nov}}[e_1,e_0].
\]
在它上面定义一个位于
$H(M;\Lambda_{0,\mathrm{nov}})[1]$ 之上的
滤过左 $L_\infty$ 模结构如下：
\begin{align}
\mathfrak n_0(e_0)&=e_1,
\tag{7.4.190.1}\label{eq:7-4-190-1}\\
\text{其余所有运算 }\mathfrak n_k&\text{ 均为零}.
\tag{7.4.190.2}\label{eq:7-4-190-2}
\end{align}
于是有滤过左 $L_\infty$ 模的正合序列
\begin{equation}
0\longrightarrow
\Lambda_{0,\mathrm{nov}}[e_1]
\longrightarrow
\Lambda_{0,\mathrm{nov}}[e_1,e_0]
\longrightarrow
\Lambda_{0,\mathrm{nov}}[e_0]
\longrightarrow0.
\tag{7.4.191}\label{eq:7-4-191}
\end{equation}

引入这个模，是为了用 $L_\infty$ 结构
重写含有格罗莫夫--威腾不变量的
式~\eqref{eq:3-8-79-4}。
下面考虑
\[
\operatorname{Hom}^{\mathrm{filt}}
\bigl(
B^{\mathrm{cyc}}(H(L;\Lambda_{0,\mathrm{nov}})[1]),
\Lambda_{0,\mathrm{nov}}
\bigr).
\]
由命题~\ref{prop:7-4-184}，
这是微分分次李代数
$\operatorname{Der}B(H(L;\Lambda_{0,\mathrm{nov}})[1])$
上的滤过微分分次模。
把 $H(M;\Lambda_{0,\mathrm{nov}})[1]$
视为平凡的
$\operatorname{Der}B(H(L;\Lambda_{0,\mathrm{nov}})[1])$-模，
也就是说其作用为零。
于是
\[
\operatorname{Hom}^{\mathrm{filt}}
\bigl(
B^{\mathrm{cyc}}(H(L;\Lambda_{0,\mathrm{nov}})[1]),
H(M;\Lambda_{0,\mathrm{nov}})[1]
\bigr)
\]
是
$\operatorname{Der}B(H(L;\Lambda_{0,\mathrm{nov}})[1])$
上的滤过左 $L_\infty$ 模。

还要指出，
\[
\begin{aligned}
&\operatorname{Hom}^{\mathrm{filt}}
\bigl(
B_0^{\mathrm{cyc}}(H(L;\Lambda_{0,\mathrm{nov}})[1]),
H(M;\Lambda_{0,\mathrm{nov}})[1]
\bigr)\\
&\qquad\cong H(M;\Lambda_{0,\mathrm{nov}})[1]
\end{aligned}
\]
是一个子模
\translatornote{源书把同构右端排作
$H(M;\Lambda_{0,\mathrm{nov}}[1])$；
但左端的同态靶已经把移位放在系数括号之外，且
$B_0^{\mathrm{cyc}}=\Lambda_{0,\mathrm{nov}}$。
故译文恢复为 $H(M;\Lambda_{0,\mathrm{nov}})[1]$。
见疑误 A-0605。}。
可以把
\[
\prod_{k>0}
\operatorname{Hom}^{\mathrm{filt}}
\bigl(
B_k^{\mathrm{cyc}}(H(L;\Lambda_{0,\mathrm{nov}})[1]),
H(M;\Lambda_{0,\mathrm{nov}})[1]
\bigr)
\]
视为商模，并将它记作
\[
\operatorname{Hom}^{\mathrm{filt}}
\bigl(
B_+^{\mathrm{cyc}}(H(L;\Lambda_{0,\mathrm{nov}})[1]),
H(M;\Lambda_{0,\mathrm{nov}})[1]
\bigr).
\]
注意，上述 $k>0$ 的直积本身并不是
相应微分分次李代数模的子模。

\begin{theorem}{7.4.192}\label{thm:7-4-192}
存在下列由有隙滤过左 $L_\infty$ 模同态组成的交换图式：

\par\medskip\noindent\begin{minipage}{\textwidth}
\centering
\resizebox{.99\textwidth}{!}{$\displaystyle
\begin{array}{ccccc}
\Lambda_{0,\mathrm{nov}}[e_1]
&\longrightarrow&
\Lambda_{0,\mathrm{nov}}[e_1,e_0]
&\longrightarrow&
\Lambda_{0,\mathrm{nov}}[e_0]
\\[-.1em]
{\scriptstyle\frac1{\ell!}GW(\,\cdot\,^{\circ},[L])}
\Big\downarrow
&&{\scriptstyle\mathfrak p^\circ}\Big\downarrow
&&\Big\downarrow
\\[-.1em]
H(M)
&\longrightarrow&
\operatorname{Hom}^{\mathrm{filt}}
\bigl(B^{\mathrm{cyc}}(H(L)),H(M)\bigr)
&\longrightarrow&
\operatorname{Hom}^{\mathrm{filt}}
\bigl(B_+^{\mathrm{cyc}}(H(L)),H(M)\bigr).
\end{array}$}
\par\smallskip
{\renewcommand{\thefigure}{7.4.10}\refstepcounter{figure}
\normalfont\bfseries 图式~\thefigure\label{diag:7-4-10}}
\end{minipage}\par\medskip

图式~\ref{diag:7-4-10} 中省略了系数环
$\Lambda_{0,\mathrm{nov}}$ 和次数移位符号 $[1]$。
其中各竖直箭头都是位于定理~\ref{thm:7-4-118} 的
\[
\mathfrak q^\circ\colon
H(M;\Lambda_{0,\mathrm{nov}})[1]
\longrightarrow
\operatorname{Der}B
\bigl(H(L;\Lambda_{0,\mathrm{nov}})[1]\bigr)
\]
之上的滤过 $L_\infty$ 模同态。
第一列中的箭头由
\[
\mathbf Y\otimes e_1
\longmapsto
\frac1{\ell!}
GW_{0,\ell+1}
\bigl(\operatorname{Op}(\mathbf Y)\otimes[L]\bigr)
\]
给出，其中
\[
\mathbf Y\in
E_\ell\bigl(H(M;\Lambda_{0,\mathrm{nov}})[1]\bigr),
\]
而 $GW_{0,\ell+1}$ 是第~\ref{subsec:3-8-9} 小节中的
亏格零格罗莫夫--威腾不变量。

还有下列交换图式；它是图式~\ref{diag:7-4-10}
的含幺版本：

\par\medskip\noindent\begin{minipage}{\textwidth}
\centering
\resizebox{.99\textwidth}{!}{$\displaystyle
\begin{array}{ccccc}
\Lambda_{0,\mathrm{nov}}[e_1]
&\longrightarrow&
\Lambda_{0,\mathrm{nov}}[e_1,e_0]
&\longrightarrow&
\Lambda_{0,\mathrm{nov}}[e_0]
\\[-.1em]
{\scriptstyle\frac1{\ell!}GW(\,\cdot\,^{\circ},[L])}
\Big\downarrow
&&{\scriptstyle\mathfrak p^\circ}\Big\downarrow
&&\Big\downarrow
\\[-.1em]
H(M)
&\longrightarrow&
\operatorname{Hom}^{\mathrm{filt}}
\bigl(B_{\mathrm{red}}^{\mathrm{cyc}}(H(L)),H(M)\bigr)
&\longrightarrow&
\operatorname{Hom}^{\mathrm{filt}}
\bigl(B_{\mathrm{red}+}^{\mathrm{cyc}}(H(L)),H(M)\bigr).
\end{array}$}
\par\smallskip
{\renewcommand{\thefigure}{7.4.11}\refstepcounter{figure}
\normalfont\bfseries 图式~\thefigure\label{diag:7-4-11}}
\end{minipage}\par\medskip

这里各竖直箭头都是位于
\[
\mathfrak q^\circ\colon
H(M;\Lambda_{0,\mathrm{nov}})[1]
\longrightarrow
\operatorname{Der}B^{\mathrm{unit}}
\bigl(H(L;\Lambda_{0,\mathrm{nov}})[1]\bigr)
\]
之上的滤过 $L_\infty$ 模同态。
图式~\ref{diag:7-4-11} 中同样省略了次数移位符号 $[1]$。
\end{theorem}

\begin{proof}
先定义
\[
\begin{aligned}
\mathfrak p^\circ_{\ell,k}\colon{}&
E_\ell\bigl(C(M;\Lambda_{0,\mathrm{nov}})[2]\bigr)
\otimes\Lambda_{0,\mathrm{nov}}[e_1,e_0]
\otimes
B_k^{\mathrm{cyc}}
\bigl(C_{1+}(L;\Lambda_{0,\mathrm{nov}})[1]\bigr)\\
&\longrightarrow S^*(M;\Lambda_{0,\mathrm{nov}}).
\end{aligned}
\]
这里的记号与命题~\ref{prop:3-8-78} 相同
\translatornote{源书把前一映射的靶写成
$S(M;\Lambda_{0,\mathrm{nov}})$，又把下列 $x$ 的定义域写成
$B_k^{\mathrm{cyc}}C(L;\Lambda_{0,\mathrm{nov}})$；
命题 3.8.78 明定 $\mathfrak p_{\ell,k}$ 取值于
$S^*(M;\Lambda_{0,\mathrm{nov}})$，且其循环 bar 输入建立在
$C(L;\Lambda_{0,\mathrm{nov}})[1]$ 上。
故译文分别补回 $*$ 与 $[1]$。见疑误 A-0606、A-0607。}。
取
\[
\mathbf Y\in
E_\ell\bigl(C(M;\Lambda_{0,\mathrm{nov}})[2]\bigr),
\qquad
x\in
B_k^{\mathrm{cyc}}
\bigl(C(L;\Lambda_{0,\mathrm{nov}})[1]\bigr)
\]
并置
\begin{equation}
\mathfrak p^\circ_{\ell,k}(\mathbf Y,e_0)(x)
=
\mathfrak p\bigl(\operatorname{Op}(\mathbf Y)\otimes x\bigr).
\tag{7.4.193}\label{eq:7-4-193}
\end{equation}
右端在第~\ref{subsec:3-8-9} 小节定义。
再置
\begin{equation}
\mathfrak p^\circ_{\ell,k}(\mathbf Y,e_1)(x)
=
\begin{cases}
\displaystyle
\frac{c}{\ell!}
GW_{0,\ell+1}\bigl(\operatorname{Op}(\mathbf Y),[L]\bigr),
&
\begin{gathered}
x=c\in
B_0^{\mathrm{cyc}}
\bigl(H(L;\Lambda_{0,\mathrm{nov}})[1]\bigr)\\
=\Lambda_{0,\mathrm{nov}},
\end{gathered}\\[1.1em]
0,&\text{其他情形}.
\end{cases}
\tag{7.4.194}\label{eq:7-4-194}
\end{equation}
式~\eqref{eq:7-4-194} 采用亏格下标 $0$
\translatornote{源书此处写 $GW_{\ell+1}$；
但定理 7.4.192 第一列及式 (3.8.79.4) 均为
$GW_{0,\ell+1}$，当前语境也只定义了这一亏格零不变量。
见疑误 A-0608。}。

可以证明，$\mathfrak p^\circ$ 是滤过左 $L_\infty$ 模同态，
当且仅当式~\eqref{eq:3-8-79-1}--\eqref{eq:3-8-79-4}
及式~\eqref{eq:3-8-10-2}--\eqref{eq:3-8-10-3} 成立。
因此，由定理~\ref{thm:3-8-9} 和命题~\ref{prop:3-8-78}，
$\mathfrak p^\circ$ 是滤过左 $L_\infty$ 模同态。
再应用定理~\ref{thm:7-4-186}，便得到所需的
滤过左 $L_\infty$ 模同态。

含幺版本可用式~\eqref{eq:3-8-79-5}--\eqref{eq:3-8-79-6}
及式 (3.8.10.4)--(3.8.10.6) 证明。
\end{proof}

\begin{theorem}{7.4.195}\label{thm:7-4-195}
设
\[
\psi\colon(M,L)\longrightarrow(M',L')
\]
是辛微分同胚。
可以选取定理~\ref{thm:7-4-192} 中的 $\mathfrak p^\circ$
\translatornote{源书此处只写 filtered left $L_\infty$
homomorphism，漏去 module；但源、靶都是上一条定理构造的模，
且该同态位于基 $L_\infty$ 同态之上。译文依型别补出“模”。
见疑误 A-0609。}，
使得存在滤过左 $L_\infty$ 模同态
\[
\begin{aligned}
\psi_*\colon{}&
\operatorname{Hom}^{\mathrm{filt}}
\bigl(
B^{\mathrm{cyc}}(H(L;\Lambda_{0,\mathrm{nov}})[1]),
H(M;\Lambda_{0,\mathrm{nov}})[1]
\bigr)\\
&\longrightarrow
\operatorname{Hom}^{\mathrm{filt}}
\bigl(
B^{\mathrm{cyc}}(H(L';\Lambda_{0,\mathrm{nov}})[1]),
H(M';\Lambda_{0,\mathrm{nov}})[1]
\bigr),
\end{aligned}
\]
它位于定理~\ref{thm:7-4-120} 的
\[
\begin{aligned}
\psi_*\colon
\operatorname{Der}B
\bigl(H(L;\Lambda_{0,\mathrm{nov}})[1]\bigr)
\longrightarrow
\operatorname{Der}B
\bigl(H(L';\Lambda_{0,\mathrm{nov}})[1]\bigr)
\end{aligned}
\]
之上，并使下图在同伦意义下交换：

\par\medskip\noindent\begin{minipage}{\textwidth}
\centering
\resizebox{.96\textwidth}{!}{$\displaystyle
\begin{array}{ccc}
\Lambda_{0,\mathrm{nov}}[e_1,e_0]
&\xrightarrow{\ \mathfrak p^\circ\ }&
\operatorname{Hom}^{\mathrm{filt}}
\bigl(B^{\mathrm{cyc}}(H(L;\Lambda_{0,\mathrm{nov}})),
H(M;\Lambda_{0,\mathrm{nov}})\bigr)
\\[-.1em]
\Big\Vert&&{\scriptstyle\psi_*}\Big\downarrow
\\[-.1em]
\Lambda_{0,\mathrm{nov}}[e_1,e_0]
&\xrightarrow{\ \mathfrak p^\circ\ }&
\operatorname{Hom}^{\mathrm{filt}}
\bigl(B^{\mathrm{cyc}}(H(L';\Lambda_{0,\mathrm{nov}})),
H(M';\Lambda_{0,\mathrm{nov}})\bigr).
\end{array}$}
\par\smallskip
{\renewcommand{\thefigure}{7.4.12}\refstepcounter{figure}
\normalfont\bfseries 图式~\thefigure\label{diag:7-4-12}}
\end{minipage}\par\medskip

含幺版本亦成立。
图式~\ref{diag:7-4-12} 中省略了次数移位符号 $[1]$。
\end{theorem}

\begin{proof}
其证明与定理~\ref{thm:7-4-120} 的证明类似，采用分叉法；故略去。
\end{proof}

\begin{remark}{7.4.196}\label{rem:7-4-196}
事实上，为了证明定理~\ref{thm:3-8-9}
和命题~\ref{prop:3-8-78}，
需要使用定理~\ref{thm:7-4-192}、\ref{thm:7-4-195}
的 $A_{n,K}$ 版本。
也就是说，先用定理~\ref{thm:3-8-9}
和命题~\ref{prop:3-8-78} 的 $A_{n,K}$ 版本，
证明定理~\ref{thm:7-4-192} 的 $A_{n,K}$ 版本；
再把它与定理~\ref{thm:7-4-195} 的 $A_{n,K}$ 版本合用，
证明定理~\ref{thm:7-4-192} 的 $A_\infty$ 版本。
后者随即蕴含定理~\ref{thm:3-8-9}
和命题~\ref{prop:3-8-78}。
\end{remark}

\begin{remark}{7.4.197}\label{rem:7-4-197}
本节中的两个运算
$\mathfrak p_{\ell-1,k}$ 与 $\mathfrak q_{\ell,k-1}$
扮演颇为不同的角色。
不过，两者的几何来源是同一个伪全纯圆盘模空间；
这些圆盘带有 $\ell$ 个内部标记点和 $k$ 个边界标记点。
区别在于把哪一个标记点（内部的或边界的）看作第 $0$ 个标记点。

在第~\ref{sec:7-2} 节和本节的构造中，
这个区别至关重要，因为选取链的可数集这一过程
破坏了标记点之间的对称性。
不过，如果采用连续扰动族的方法（见下一节），
就可以保持这种对称性；见 \textbf{[Fuk05II]}。
因此，对
\[
P_i\in C(L;\mathbb R),
\qquad
Q_i\in C(M;\mathbb R),
\]
有
\begin{equation}
\begin{aligned}
&\left\langle
P_1,
\mathfrak q_{\ell,k-1}
(Q_1,\ldots,Q_\ell;P_2,\ldots,P_k)
\right\rangle_L\\
&\qquad=
\pm\left\langle
Q_1,
\mathfrak p_{\ell-1,k}
(Q_2,\ldots,Q_\ell;P_1,\ldots,P_k)
\right\rangle_M.
\end{aligned}
\tag{7.4.198}\label{eq:7-4-198}
\end{equation}
式~\eqref{eq:7-4-198} 左端列至 $P_k$
\translatornote{源书左端只列 $P_2,\ldots,P_{k-1}$；
但 $\mathfrak q_{\ell,k-1}$ 需要 $k-1$ 个边界输入，
而紧接的式 (7.4.199) 也明确列出 $P_2,\ldots,P_k$。
故译文补入 $P_k$。见疑误 A-0610。}。
这里 $\langle\ ,\ \rangle_L$、
$\langle\ ,\ \rangle_M$
分别是 $L$、$M$ 的庞加莱对偶配对。

式~\eqref{eq:7-4-198} 蕴含猜想~\ref{conj:3-6-48}
的下列推广：
\begin{equation}
\begin{aligned}
&\left\langle
P_1,
\mathfrak q_{\ell,k-1}
(Q_1,\ldots,Q_\ell;P_2,\ldots,P_k)
\right\rangle_L\\
&\qquad=
\pm\left\langle
P_k,
\mathfrak q_{\ell,k-1}
(Q_1,\ldots,Q_\ell;P_1,P_2,\ldots,P_{k-1})
\right\rangle_L.
\end{aligned}
\tag{7.4.199}\label{eq:7-4-199}
\end{equation}
利用连续扰动族的方法，也可以在 $\mathbb R$ 系数上证明
式~\eqref{eq:7-4-199}。
作者不知道如何在 $\mathbb Q$ 系数上证明
式~\eqref{eq:7-4-198}、\eqref{eq:7-4-199}。

可以利用式~\eqref{eq:7-4-198}、\eqref{eq:7-4-199}
在 Hochschild 复形和循环 bar 复形上引入更多结构。
这些结构可用于研究弗洛尔理论的高亏格推广。
正如注~\ref{rem:7-4-42}~\textup{(2)} 中已经提到的，
这也与弦拓扑 \textbf{[CS99]} 有关。
作者计划在不久的将来另文系统讨论这些结构。

T. Tradler 在 \textbf{[Tra02]} 中讨论了把“同伦内积”
与 $A_\infty$ 结构作为彼此分立的运算来研究的想法。
第~\ref{subsec:7-4-10}、\ref{subsec:7-4-11} 小节
对循环 bar 复形以及同态 $\mathfrak p,\mathfrak q$ 的讨论，
似乎与 Tradler 的研究有关。
\end{remark}
